\section{黎曼度量让弯曲空间有长度}
\label{sec:12-Real-Analysis-Support/08-Manifold-Basics/03}

北京飞纽约的航线在机上屏幕的世界地图里向北拱出一个大弯，几乎蹭到北极圈。第一次看见的人多半以为航司在绕路。把同样两点连成地图上的一条直线，量出来的长度反而更大——那条直线在球面上不是最短的走法。地图没画错，飞机也没绕远，两者只是用了两把不同的尺：一把是屏幕上的平面尺，另一把是地球表面自己的尺。

\begin{center}
\begin{tikzpicture}[>=stealth,scale=1.0]
% 平面世界地图
\draw[gray!55] (0,0) rectangle (9,4.2);
\foreach \y in {1.0,2.1,3.2} \draw[gray!22] (0,\y) -- (9,\y);
\foreach \x in {1.5,3.0,4.5,6.0,7.5} \draw[gray!22] (\x,0) -- (\x,4.2);
\coordinate (B) at (7.2,2.6);
\coordinate (N) at (1.9,2.5);
\fill (B) circle (2pt); \node[below=2pt] at (B) {北京};
\fill (N) circle (2pt); \node[below=2pt] at (N) {纽约};
% 地图上的直线
\draw[gray!70!black,dashed,line width=1pt] (N) -- (B);
\node[gray!70!black,font=\small] at (4.55,2.32) {图上的直线};
% 大圆在地图上的投影
\draw[orange!85!black,line width=1.3pt]
  plot[domain=0:1,samples=60]
  ({1.9+5.3*\x},{2.55+1.15*sin(180*\x)});
\node[orange!85!black,font=\small] at (4.55,4.0) {球面上的最短路};
\end{tikzpicture}
\end{center}

\emph{同一对端点，同一张图，两条曲线的长短之争取决于用哪把尺去量。}

\subsection{度量是逐点装在切空间上的一个内积}

上一节把方向找回来了：每点一个 \(d\) 维切空间。可是向量空间的公理里没有长度这回事，\(T_{p}M\) 中一个向量有多长、两个向量夹角多大，目前无从回答。补法只有一个——逐点指定内积。

\begin{definition}[黎曼度量]
流形 \(M\) 上的黎曼度量 \(g\) 为每个 \(p\in M\) 指定一个对称双线性型 \(g_{p}:T_{p}M\times T_{p}M\to\R\)，满足 \(g_{p}(v,v)>0\)（\(v\neq0\)），并且在任一坐标卡下，分量 \(g_{ij}(p)=g_{p}(\partial_{i},\partial_{j})\) 随 \(p\) 光滑变化。配了 \(g\) 的流形叫黎曼流形。
\end{definition}

最省事的例子是诱导度量：\(M\subseteq\R^{n}\) 是子流形时，\(T_{p}M\) 本身就是 \(\R^{n}\) 的子空间，把环境的欧氏内积限制上去即可。球面、环面、数据流形上默认的那把尺都是这么来的。

有了 \(g\)，曲线 \(\gamma:[a,b]\to M\) 的长度写成

\[
L(\gamma)=\int_{a}^{b}\sqrt{g_{\gamma(t)}\bigl(\gamma'(t),\gamma'(t)\bigr)}\,dt,
\]

两点间的距离定义为所有连接它们的曲线长度的下确界。夹角、正交、Cauchy--Schwarz 不等式在每个 \(T_{p}M\) 内部原样可用，因为 \(g_{p}\) 就是一个普通内积。要检查这套定义配不配得上``距离''这个名字，得核对三角不等式与非退化性；前者靠路径拼接，后者靠 \(g\) 的正定性加上流形的 Hausdorff 性质，两处都不难，但都不是白送的。

\subsection{有了度量，``直线''的位置被测地线接管}

欧氏空间里直线实现最短。黎曼流形上没有直线可用，接班的是测地线：在诱导度量下，它是加速度始终垂直于流形的曲线——沿途不主动拐弯，弯的部分全由流形逼着弯。

球面上的测地线是大圆，也就是过球心的平面与球面的交线。开头那条航线就是大圆在平面地图上的投影，之所以看着朝北鼓出去，是因为墨卡托投影把高纬度横向拉宽了，图上等长的一段在高纬处对应球面上更短的一段。

有三处必须分清。距离是内在的：大圆的长度用球面自己的度量算出来，全程没有离开球面。直线概念被替换掉了，测地线不是任何坐标下的直线，坐标直线在球面上一般不测地。最短还只是局部性质：沿大圆走过对径点之后再往前，仍是测地线，却早已不是最短路；南北极之间的经线有无穷多条，条条一样短。测地线保证局部最短，全局最短既可能不唯一，也可能压根不存在。

\subsection{同一个流形配不同的度量，几何完全不同}

到这里该点破一件事：\textbf{度量是外加的结构，不是流形自带的。} 流形的定义里只有卡和转移映射，没有一个字提到长度。\(g\) 是后来放上去的，而且可以有很多种放法。

拿 \(\R^{2}\) 举例。装上标准欧氏内积，测地线是直线，三角形内角和是 \(\pi\)。改成 \(g_{p}(v,w)=v^{\trans}A^{\trans}Aw\)（\(A\) 可逆常矩阵），得到的仍是平坦几何，只是把``单位圆''换成了椭圆——马氏距离就是这么回事。再看上半平面 \(\set{(x,y):y>0}\)，装上 \(g_{(x,y)}(v,w)=\ip{v}{w}/y^{2}\)，同一个流形立刻变成双曲平面：测地线成了垂直于 \(x\) 轴的半圆，三角形内角和小于 \(\pi\)，从一点出发的圆周长随半径指数增长。底下那个可微结构一动没动，几何却换了一副面孔。

这解释了为什么本节标题里是``让''而不是``发现''。长度不是从流形里挖出来的，是装上去的。装哪一把尺，是建模者的决定。

\subsection{Fisher 度量是统计流形上的一种装法}

参数族 \(\set{p_{\theta}:\theta\in\Theta}\) 是这条原则最有说服力的场合。参数向量住在 \(\R^{k}\) 里，欧氏距离随手可得，可是它量的是参数数值的差，跟``两个分布有多容易区分''没有必然联系：一个方差参数挪动 \(0.01\)，在样本量大时可能天翻地覆，另一个几乎无关紧要的参数挪动 \(1\) 也许什么都不改变。

Fisher 信息提供了另一把尺：

\[
g_{\theta}=\E_{\theta}\!\left[
\nabla_{\theta}\log p_{\theta}\,
\bigl(\nabla_{\theta}\log p_{\theta}\bigr)^{\trans}
\right].
\]

它在 \(\theta\) 处按对数似然的波动来量方向，方向上分布变化越剧烈，度量给出的长度越大。Part 11 已经用它把微小参数位移换算成 KL 散度的二阶项；升级为黎曼度量之后，参数空间成了统计流形，而自然梯度不过是在这个几何里重新算的最速下降方向——梯度下降沿欧氏最陡方向走，自然梯度沿 Fisher 意义下最陡的方向走，两者在同一个参数空间上给出不同的轨迹。\hyperref[ch:105]{``信息几何：参数空间中的分布几何''}整章都在展开这条线。

\subsection{换一把尺，下游结论跟着换}

这件事对做数据的人有直接后果。相似度、聚类、近邻检索、插值路径，每一个都建立在某个距离之上，而那个距离是选出来的，不是数据发的。

Isomap 是个干净的例子。它主张两点之间的距离应当沿数据流形的表面量，于是用邻接图上的最短路去近似测地距离。瑞士卷上下相邻两层在环境空间里挨得很近，沿曲面却隔着大半圈，这个差别正是诱导度量与环境度量的分歧。子流形上两套距离始终并存，关系有两条：内在距离不小于环境距离，因为直线不会比任何路径更长；两点足够近时二者一阶吻合，因为切空间是一阶近似。Isomap 的赌注全押在这两条上——短边可以直接用环境距离，长距离必须沿图走。噪声和稀疏采样弄坏的，恰是``足够近''与``足够远''之间那条线。

再往前一步的说法要克制。度量的选择改变结论，这是数学事实；某个具体任务上该选哪把尺，几何不给答案，那属于建模判断，要靠对数据生成过程的了解和下游指标去定。能确定的只是：如果结论在换尺之后翻转，那么它依赖的是尺，不是数据。

\subsection{三层地基铺完了}

可拉平（流形）、有方向（切空间）、有长度（黎曼度量），三层叠起来，弯曲空间上的微积分才有落脚处。接下来该把这副地基用到对象本身：下一节\hyperref[sec:12-Real-Analysis-Support/08-Manifold-Basics/04]{``流形上的函数与映射''}讨论怎样对流形上的函数求导，以及为什么链式法则在这里一个字都不用改。
